\clearpage
\phantomsection
\chapter*{Introducción}
\addcontentsline{toc}{chapter}{Introducción}

El procesamiento de lenguaje natural depende de la disponibilidad de datos y recursos lingüísticos adecuados. En Quechua Collao, los trabajos disponibles incluyen corpus paralelo y recursos léxicos que aportan materiales para el análisis computacional, aunque no constituyen por sí mismos un asistente de generación local basado en recuperación documental \parencite{Orcotoma2025,Duran2026}.

La investigación reciente también muestra que el uso de contexto en lenguas de bajos recursos depende de la pertinencia y del tipo de información suministrada. En tareas de traducción de quechua sureño al español, \textcite{Court2024} observaron resultados distintos según la información contextual recuperada; el estudio no evalúa generación en Quechua Collao. En un dominio educativo diferente, \textcite{Hevia2025} analizaron un sistema de tutoría con un modelo pequeño y RAG offline, y señalaron dificultades cuando se incorporan contextos extensos o irrelevantes. Estos resultados orientan el diseño de la evaluación, pero no permiten anticipar el rendimiento de la arquitectura propuesta.

La ejecución local puede reducir la dependencia de servicios remotos, aunque debe ajustarse a los límites de memoria y cómputo del dispositivo. Las revisiones sobre modelos en dispositivos describen la cuantización y otras técnicas de compresión como opciones para reducir recursos, junto con compromisos de rendimiento que deben medirse \parencite{Xu2024}. La familia Qwen2.5 ofrece modelos de pesos abiertos en varios tamaños; el proyecto considera como candidato un modelo compacto de esa familia, cuya elección final dependerá de las pruebas de compatibilidad y memoria \parencite{qwen2024}.

QLoRA adapta un modelo cuantizado mediante matrices de bajo rango y reduce el costo de ajuste respecto del entrenamiento completo en las condiciones evaluadas por sus autores \parencite{Dettmers2023}. Sin embargo, ese antecedente no demuestra que el ajuste de un modelo de tres mil millones de parámetros sea viable en la GPU local de 4 GB disponible para este proyecto. Por ello, la factibilidad del ajuste se tratará como una condición experimental y, si los recursos no bastan, se documentará el alcance de la alternativa evaluada.

La revisión preliminar identifica estudios sobre recursos quechuas, uso de contexto para traducción y sistemas de tutoría offline, pero todavía no permite afirmar de manera concluyente que no existan trabajos comparables. La brecha de investigación se formula provisionalmente como la falta de evidencia reunida en este proyecto sobre el comportamiento de una arquitectura que combine un modelo compacto, recuperación documental local y adaptación QLoRA para generar respuestas contextualizadas en Quechua Collao sin conexión.

En consecuencia, la investigación pregunta de qué manera el diseño e implementación de una arquitectura SLM--RAG con QLoRA influye en la generación contextualizada en Quechua Collao dentro de un entorno computacional offline. El plan presenta el problema, los objetivos y las hipótesis; después desarrolla el estado del arte y el marco teórico, la metodología, la factibilidad y los aspectos administrativos. Los resultados, la discusión y las conclusiones corresponderán al informe final una vez realizadas las pruebas.
